\subsection{Time Lax matrix from the finite-lattice construction}
\label{sec:class6-time-lax}

We apply the finite-lattice construction of
Eq.~\eqref{eq:finite-lattice-time-Lax} to the Class 6 operators of
Section~\ref{sec:class6-spatial-lax}.  As in the Class 5 calculation, we
remove a field-independent scalar term by choosing
\begin{equation}
\widehat A_{6;a,n}
:=A_{6;a,n}+\frac{\ii}{u}\Id,
\label{eq:class6-normalized-time-lax}
\end{equation}
where the identity acts on the auxiliary space and the physical sites
$n-1,n$.  The normalized spatial operator then gives
\begin{equation}
\widehat A_{6;a,n}
=-\ii\widehat L_{6;a,n}^{-1}
\left(
[h_{6;n-1,n},\widehat L_{6;a,n}]
+\partial_u\widehat L_{6;a,n}
\right).
\label{eq:class6-time-lax-from-Lhat}
\end{equation}
The derivative is taken at fixed lattice coupling $\kappa_6$.  With the
scaling in Eq.~\eqref{eq:class6-continuum-scaling}, the exact expansion
\eqref{eq:class6-L-expansion} therefore implies
\begin{equation}
\partial_u\widehat L_{6;a,n}
=\epsilon^2\partial_\lambda L^{(1)}_{6;a,n}
+\epsilon^3\partial_\lambda L^{(2)}_{6;a,n}
+\epsilon^4\partial_\lambda L^{(3)}_{6;a,n}.
\label{eq:class6-u-derivative-Lhat}
\end{equation}
Here and below, $\partial_\lambda$ holds $\alpha_6$ fixed, and also holds the
spin fixed when acting on a classical matrix.

The bond Hamiltonian is
\begin{equation}
h_{6;n-1,n}
=h^{(0)}_{n-1,n}+\epsilon^2\alpha_6K_{6;n-1,n},
\qquad h^{(0)}:=2P.
\label{eq:class6-h-expansion}
\end{equation}
Expanding Eq.~\eqref{eq:class6-time-lax-from-Lhat} gives
\begin{equation}
\widehat A_{6;a,n}
=\epsilon A^{(1)}_{6;a,n}
+\epsilon^2 A^{(2)}_{6;a,n}
+\mathcal O(\epsilon^3),
\label{eq:class6-A-expansion}
\end{equation}
where
\begin{align}
A^{(1)}_{6;a,n}
&=-\ii[h^{(0)}_{n-1,n},L^{(1)}_{6;a,n}],
\label{eq:class6-A1}\\
A^{(2)}_{6;a,n}
&=\ii L^{(1)}_{6;a,n}
[h^{(0)}_{n-1,n},L^{(1)}_{6;a,n}]
\nonumber\\
&\quad-\ii[h^{(0)}_{n-1,n},L^{(2)}_{6;a,n}]
-\ii\partial_\lambda L^{(1)}_{6;a,n}.
\label{eq:class6-A2}
\end{align}
The deformation of the bond Hamiltonian first contributes explicitly at
order $\epsilon^3$, since its commutator with the leading identity in
$\widehat L_6$ vanishes.  The dependence on $\alpha_6$ in
Eqs.~\eqref{eq:class6-A1} and \eqref{eq:class6-A2} enters through
$L_6^{(1)}$ and $L_6^{(2)}$.

Taking the lower symbol on the two physical sites gives
\begin{equation}
\bigl(A^{(1)}_{6;a,n}\bigr)^\downarrow
=-2(\boldsymbol S_{n-1}\times\boldsymbol S_n)^i
\frac{\partial U_6}{\partial S^i}.
\label{eq:class6-A1-lower-symbol}
\end{equation}
The Cartesian derivatives in this formula act on the explicit matrix
\eqref{eq:class6-spatial-Lax}, with $S^\pm=S^1\pm\ii S^2$, before imposing
the unit-spin constraint.  Since $U_6$ is affine in the spin components,
these derivatives are independent of the spin.  At $x=x_n$,
\begin{equation}
\boldsymbol S_{n-1}\times\boldsymbol S_n
=\epsilon\boldsymbol S\times\partial_x\boldsymbol S
+\mathcal O(\epsilon^2).
\label{eq:class6-neighbour-cross-product}
\end{equation}
The second time coefficient is needed only at coincident spin arguments.
Using $S^+S^-+(S^3)^2=1$ and the traceless spatial matrix
$U_{6,0}$ of Eq.~\eqref{eq:class6-traceless-U}, we obtain
\begin{equation}
\left.\bigl(A^{(2)}_{6;a,n}\bigr)^\downarrow
\right|_{\boldsymbol S_{n-1}=\boldsymbol S_n=\boldsymbol S}
=-2\ii\partial_\lambda U_{6,0}
+\ii\alpha_6S^+S^3\Id_{\mathcal V_a}.
\label{eq:class6-A2-lower-symbol}
\end{equation}
Consequently,
\begin{equation}
\widehat A_{6;a,n}^\downarrow
=\epsilon^2V_6^\downarrow(x_n,t;\lambda)
+\mathcal O(\epsilon^3),
\label{eq:class6-A-lower-expansion}
\end{equation}
with
\begin{equation}
V_6^\downarrow
=-2(\boldsymbol S\times\partial_x\boldsymbol S)^i
\frac{\partial U_6}{\partial S^i}
-2\ii\partial_\lambda U_{6,0}
+\ii\alpha_6S^+S^3\Id_{\mathcal V_a}.
\label{eq:class6-Vdown-direct}
\end{equation}

The products in the discrete zero-curvature equation require a separate
calculation.  Let $\widehat{\mathcal D}^{(\pm)}_{6;a,n}$ denote
Eqs.~\eqref{eq:Rplus-hat-definition} and
\eqref{eq:Rminus-hat-definition} evaluated for Class 6.  For the first
product, the operator expansions give
\begin{align}
\widehat{\mathcal D}^{(+)}_{6;a,n}
={}&\epsilon^2\Bigl[
\bigl(A^{(1)}_{6;a,n+1}L^{(1)}_{6;a,n}\bigr)^\downarrow
-\bigl(A^{(1)}_{6;a,n+1}\bigr)^\downarrow
 \bigl(L^{(1)}_{6;a,n}\bigr)^\downarrow
\Bigr]
\nonumber\\
&+\epsilon^3\Bigl[
\bigl(A^{(2)}_{6;a,n+1}L^{(1)}_{6;a,n}
+A^{(1)}_{6;a,n+1}L^{(2)}_{6;a,n}\bigr)^\downarrow
\nonumber\\
&\hspace{3.4em}
-\bigl(A^{(2)}_{6;a,n+1}\bigr)^\downarrow
 \bigl(L^{(1)}_{6;a,n}\bigr)^\downarrow
-\bigl(A^{(1)}_{6;a,n+1}\bigr)^\downarrow
 \bigl(L^{(2)}_{6;a,n}\bigr)^\downarrow
\Bigr]
+\mathcal O(\epsilon^4).
\label{eq:class6-overlapping-product-expansion}
\end{align}
The expression for $\widehat{\mathcal D}^{(-)}_{6;a,n}$ has the spatial
factor on the left and the time factor at $n$, supported on $n-1,n$.
Terms involving the third time coefficient cancel within each difference,
because they multiply the leading identity in $\widehat L_6$.
The coefficient $L_6^{(3)}$ first enters these differences at order
$\epsilon^4$.  The continuum source therefore requires the products
displayed in Eq.~\eqref{eq:class6-overlapping-product-expansion}, including
those involving $L_6^{(2)}$.

At coincident spins, the order-$\epsilon^2$ terms cancel between
$\widehat{\mathcal D}^{(+)}_{6;a,n}$ and
$\widehat{\mathcal D}^{(-)}_{6;a,n}$.  To express the surviving contribution,
define $\Xi_6$ as the difference between the lower symbol of the square of
the first spatial coefficient and the square of its lower symbol:
\begin{align}
\Xi_6
&:=\left[\bigl(L_6^{(1)}\bigr)^2\right]^\downarrow-U_6^2
\nonumber\\
&=\frac{1}{4\lambda^2}\Id_{\mathcal V_a}
+2\bigl(L_6^{(2)}\bigr)^\downarrow-U_6^2.
\label{eq:class6-symbol-square-difference}
\end{align}
The second equality follows from Eq.~\eqref{eq:class6-L-first-square}.
In particular, Eq.~\eqref{eq:class6-K-lower-symbols} gives
\begin{equation}
\bigl(L_6^{(2)}\bigr)^\downarrow
=\begin{pmatrix}
\alpha_6 S^+/4&\alpha_6 S^3/2-\alpha_6^2\lambda^2S^+\\
0&-\alpha_6 S^+/4
\end{pmatrix}.
\label{eq:class6-second-spatial-lower-symbol}
\end{equation}
Expanding the spin arguments of the order-$\epsilon^2$ terms to first
order, and evaluating the order-$\epsilon^3$ products at coincident spins,
yields
\begin{equation}
\mathcal C_6
:=\lim_{\epsilon\to0}\epsilon^{-3}
\left(\widehat{\mathcal D}^{(+)}_{6;a,n}
-\widehat{\mathcal D}^{(-)}_{6;a,n}\right)
=2\ii\bigl(\partial_x\Xi_6+[\Xi_6,U_6]\bigr).
\label{eq:class6-shared-site-correction}
\end{equation}
This calculation uses the unit-spin constraint and its spatial derivative
and holds for arbitrary time dependence of the spin.

The local time correction
\begin{align}
\Delta V_6
&:=2\ii\Xi_6-\frac{\ii}{4\lambda^2}\Id_{\mathcal V_a}
\nonumber\\
&=4\ii\bigl(L_6^{(2)}\bigr)^\downarrow
-2\ii U_6^2+\frac{\ii}{4\lambda^2}\Id_{\mathcal V_a}
\label{eq:class6-DeltaV}
\end{align}
then satisfies
\begin{equation}
\partial_x\Delta V_6+[\Delta V_6,U_6]=\mathcal C_6.
\label{eq:class6-DeltaV-absorbs-C}
\end{equation}
The field-independent scalar in Eq.~\eqref{eq:class6-DeltaV} fixes a
representative with traceless final time matrix.  To simplify it, the
explicit spatial matrix gives
\begin{align}
\det U_6&=-\frac{\alpha_6}{2}S^+S^3,
\nonumber\\
\partial_\lambda U_{6,0}
&=4\bigl(L_6^{(2)}\bigr)^\downarrow-\frac{1}{\lambda}U_{6,0}.
\label{eq:class6-spatial-spectral-identities}
\end{align}
The determinant identity uses the unit-spin constraint.  The
Cayley--Hamilton identity, together with $\tr U_6=1/(2\lambda)$, now
reduces Eq.~\eqref{eq:class6-DeltaV} to
\begin{equation}
\Delta V_6
=\ii\partial_\lambda U_{6,0}
-\ii\alpha_6S^+S^3\Id_{\mathcal V_a}.
\label{eq:class6-DeltaV-spectral}
\end{equation}
Adding this correction to Eq.~\eqref{eq:class6-Vdown-direct} yields
\begin{equation}
V_6:=V_6^\downarrow+\Delta V_6
=-2(\boldsymbol S\times\partial_x\boldsymbol S)^i
\frac{\partial U_6}{\partial S^i}
-\ii\partial_\lambda U_{6,0},
\qquad \tr V_6=0.
\label{eq:class6-time-Lax-compact}
\end{equation}

We write $V_6(x,t;\lambda)=\bigl((V_6)_{rs}\bigr)_{r,s=1}^{2}$ and
$S_x^a:=\partial_x S^a$ for $a\in\{+,-,3\}$.  The matrix elements are
\begin{align}
(V_6)_{11}
={}&\frac{\ii}{4\lambda^2}\Bigl[
S^3-2\alpha_6\lambda^2S^+
+\lambda(S^-S_x^+-S^+S_x^-)
\nonumber\\
&\hspace{3.2em}
+4\alpha_6\lambda^3(S^+S_x^3-S^3S_x^+)
\Bigr],
\label{eq:class6-time-Lax-11}\\
(V_6)_{12}
={}&\frac{\ii}{4\lambda^2}\Bigl[
S^--4\alpha_6\lambda^2S^3+12\alpha_6^2\lambda^4S^+
\nonumber\\
&\hspace{3.2em}
+2\lambda(S^3S_x^--S^-S_x^3)
+4\alpha_6\lambda^3(S^-S_x^+-S^+S_x^-)
\nonumber\\
&\hspace{3.2em}
+8\alpha_6^2\lambda^5(S^3S_x^+-S^+S_x^3)
\Bigr],
\label{eq:class6-time-Lax-12}\\
(V_6)_{21}
={}&\frac{\ii}{4\lambda^2}\Bigl[
S^++2\lambda(S^+S_x^3-S^3S_x^+)
\Bigr],
\label{eq:class6-time-Lax-21}\\
(V_6)_{22}
={}&-(V_6)_{11}.
\label{eq:class6-time-Lax-22}
\end{align}

The role of the second spatial coefficient is also visible directly in
the source.  Expanding Eq.~\eqref{eq:class6-shared-site-correction} gives
\begin{equation}
\mathcal C_6
=-2\ii\partial_x(U_6^2)
+4\ii\partial_x\bigl(L_6^{(2)}\bigr)^\downarrow
+4\ii\left[\bigl(L_6^{(2)}\bigr)^\downarrow,U_6\right].
\label{eq:class6-source-expanded}
\end{equation}
At a fixed time, consider the spatially uniform spin configuration
$\boldsymbol S=(0,0,1)^{\mathsf T}$.  Then
\begin{equation}
U_6=\begin{pmatrix}
1/(2\lambda)&\alpha_6\lambda\\
0&0
\end{pmatrix},
\qquad
\mathcal C_6=\begin{pmatrix}
0&-\ii\alpha_6/\lambda\\
0&0
\end{pmatrix}.
\label{eq:class6-uniform-source}
\end{equation}
The source is nonzero for $\alpha_6\ne0$.
Omitting the term involving $L_6^{(2)}$ from
Eq.~\eqref{eq:class6-DeltaV} would leave a polynomial in $U_6$, whose
spatial derivative and commutator with $U_6$ both vanish in this example.
The total-derivative correction sufficient for Class 5 therefore fails
for Class 6.

In Section~\ref{sec:class6-eom} we derive the Hamiltonian equations
independently and establish their off-shell equivalence to the
zero-curvature equation for $(U_6,V_6)$.
